\section{Layers}\label{sec:layers}

The later sections use Theorem~\ref{thm:forest} in two ways: by Corollary~\ref{cor:cycles}, and
by the lemma of this section.

\begin{idea}{the content of the lemma}
If the outside part is a forest for more than half of the even subgraphs, by a fixed margin, then
every distance layer around the marked set is small. In the other case a window of consecutive layers is a large set
with a small boundary, and the forest estimate applies inside this set. We use the lemma for
paths in Section~\ref{sec:poles} and for pieces in Section~\ref{sec:half}.
\end{idea}

Let $G$ be a multigraph and $P$ a set of vertices. A \emph{level function} is a map
$\lambda\colon V(G)\to\mathbb{Z}$ such that $|\lambda(u)-\lambda(w)|\le1$ for every edge $uw$ that
has an end outside $P$. The distance from a fixed vertex is a level function. For
$\delta\in(0,\half)$ let $R_\delta\ge2^{100}$ be an integer such that every $r\ge R_\delta$ is in
the range of Theorem~\ref{thm:forest} and satisfies $20/\lgk3r\le\delta$. For an integer $k\ge1$
put
\[
  \Theta_\delta(k)=\max\bigl\{2R_\delta,\;2^{100}(k+1)^{2}\bigr\}.
\]

\begin{lemma}[Layers]\label{lem:width}
Let $G$ be a connected subcubic multigraph, and let $P$ be a nonempty set of vertices such that
$G[P]$ is connected and every vertex outside $P$ has degree three. Let $\lambda$ be a level
function such that at most $k$ vertices of $P$ have any given level. Let $\delta\in(0,\half)$, and
suppose that $q(G,P)>\half+\delta$. Then, for every $j$, fewer than $\Theta_\delta(k)$ vertices
outside $P$ have level $j$.
\end{lemma}

\begin{figure}[ht]
\centering
\begin{tikzpicture}[x=1cm,y=1cm]
  % the graph G
  \draw[draw=slate!70, fill=black!4, line width=0.6pt] (0,0) ellipse [x radius=1.25, y radius=1.15];
  \draw[draw=slate!70, fill=black!12, line width=0.6pt] (-0.4,0.15) ellipse [x radius=0.45, y radius=0.4];
  \node[text=slate] at (-0.4,0.15) {\small$P$};
  \node[note] at (0.15,-0.75) {$C_0$};
  \draw[region] (4.0,0) ellipse [x radius=1.45, y radius=1.15];
  \draw[accent, line width=0.5pt, dash pattern=on 2pt off 2pt] (4.4,0.1) ellipse [x radius=0.75, y radius=0.6];
  \node[text=accent] at (4.4,0.1) {$X$};
  \node[note, text=accent] at (4.0,-0.85) {$X^{+}$};
  \foreach \y in {0.6,0.2,-0.2,-0.6}{
    \draw[edge] (1.05,\y) -- (2.8,\y);
    \node[vtx] at (1.05,\y) {};
    \node[vtx] at (2.8,\y) {};}
  \node[note] at (1.95,-1.6) {the graph $G$};
  % the multigraph Gamma
  \begin{scope}[xshift=8.3cm]
    \draw[hot] (0.5,0) ellipse [x radius=0.33, y radius=0.85];
    \draw[region] (3.7,0) ellipse [x radius=1.45, y radius=1.15];
    \draw[accent, line width=0.5pt, dash pattern=on 2pt off 2pt] (4.1,0.1) ellipse [x radius=0.75, y radius=0.6];
    \node[text=accent] at (4.1,0.1) {$X$};
    \node[note, text=accent] at (3.7,-0.85) {$X^{+}$};
    \foreach \y/\x in {0.6/0.73,0.2/0.82,-0.2/0.82,-0.6/0.73}{
      \draw[edge] (\x,\y) -- (2.5,\y);
      \node[vtx] at (\x,\y) {};
      \node[vtx] at (2.5,\y) {};}
    \node[note, text=warm] at (0.5,1.15) {$W$};
    \node[note, text=warm] at (0.55,-1.2) {$P^{+}$};
    \node[note] at (1.9,-1.6) {the multigraph $\Gamma$};
  \end{scope}
\end{tikzpicture}
\caption{Step~1 in the proof of Lemma~\ref{lem:width}. In $G$ the set $X^{+}$ is joined by $d^{+}$ edges to the
connected rest $C_0$, which contains $P$. In $\Gamma$ a cycle $W$ through $d^{+}$ new vertices
takes the place of $C_0$. The even subgraphs of $G$ and of $\Gamma$ give the same distribution
of the edges inside $X^{+}$.}
\label{fig:local}
\end{figure}

\begin{proof}
Put $\Phi(m)=m/\log^{2}m$ and $\Theta=\Theta_\delta(k)$.

\step{Step 1: no large set outside $P$ has a small boundary.} We claim that no set
$X\subseteq V(G)\sm P$ has $|X|=m\ge\Theta$ and $|\delta X|\le\Phi(m)/4$. Suppose that $X$ is such a
set. Let $C_0$ be the component of $G-X$ that contains $P$, and put $X^{+}=V(G)\sm V(C_0)$. Then
$X\subseteq X^{+}$, the set $X^{+}$ avoids $P$, and $G-X^{+}=C_0$ is connected. Different components
of $G-X$ have no edge between them, so $\delta X^{+}\subseteq\delta X$. Every vertex of $X^{+}$ has
degree three. Put $d^{+}=|\delta X^{+}|\ge1$.

We form a multigraph $\Gamma$ from $G[X^{+}]$ (Figure~\ref{fig:local}). For each edge of $\delta X^{+}$ add a new vertex
and join it to the end of that edge in $X^{+}$. Join the $d^{+}$ new vertices by a cycle $W$
through all of them, and let $P^{+}$ be the set of the new vertices. Every component of
$G[X^{+}]$ has an edge to $C_0$, so $\Gamma$ is connected. It is subcubic, $\Gamma[P^{+}]$ is
connected, the vertices outside $P^{+}$ have degree three, and $\Gamma-P^{+}=G[X^{+}]$. The
multigraph $\Gamma$ has $d^{+}$ more vertices and $2d^{+}$ more edges than $G[X^{+}]$, and
$3|X^{+}|=2|E(G[X^{+}])|+d^{+}$. So its rank is
\[
  \rho=\frac{|X^{+}|+d^{+}}2+1\;\ge\;\frac m2\;\ge\;R_\delta .
\]
The function $\Phi$ is increasing on $[8,\infty)$, and $\Phi(m)/4\le\Phi(m/2)$. So
$d^{+}\le\Phi(m)/4\le\Phi(m/2)\le\Phi(\rho)$, and Theorem~\ref{thm:forest} with the choice of
$R_\delta$ gives $q(\Gamma,P^{+})\le\half+\delta$.

The set $\xi\cap E(G[X^{+}])$ has the same distribution for $\xi$ uniform on $\Z(G)$ and for
$\xi$ uniform on $\Z(\Gamma)$. Indeed, contract the connected set $V(C_0)$ of $G$ to one vertex,
and delete the loops that arise. Call the result $G'$. An even subgraph of $G$ restricts to an
even subgraph of $G'$. Every even subgraph of $G'$ arises in this way: the missing edges inside
the contracted set can be chosen in a spanning tree of that set. So the restriction is a
surjective linear map $\Z(G)\to\Z(G')$, and it sends the uniform distribution to the uniform
distribution. The contraction of $W$ in $\Gamma$ gives the same multigraph $G'$, with the new
edges in the place of the edges of $\delta X^{+}$. Hence
\[
  \PP\bigl(\xi\cap E(G[X^{+}])\text{ is a forest}\bigr)\;=\;q(\Gamma,P^{+})\;\le\;\frac12+\delta .
\]
But $G[X^{+}]$ is a subgraph of $G-P$, so this probability is at least $q(G,P)>\half+\delta$. This
contradiction proves the claim.

\step{Step 2: a window of levels.} Let $m_j$ be the number of vertices outside $P$ of level $j$.
For $i\le i'$ let $X(i,i')$ be the set of vertices outside $P$ with a level in $[i,i']$. An edge
that leaves $X(i,i')$ ends outside $P$ at level $i-1$ or $i'+1$, or it ends in $P$ at a level in
$[i-1,i'+1]$. So
\begin{equation}\label{eq:window}
  |\delta X(i,i')|\;\le\;3m_{i-1}+3m_{i'+1}+3k(i'-i+3).
\end{equation}
Suppose that $m_j\ge\Theta$ for some $j$. Start with the window $i=i'=j$. As long as a level next
to the window has more than $m^{3/4}$ vertices outside $P$, where $m=|X(i,i')|$ is the present
size, add that level to the window. The graph is finite, so this stops. A step increases
$m^{1/4}$ by more than $\frac16$, because $(t+\frac16)^{4}\le t^{4}+t^{3}$ for $t\ge1$. Let $\mu$
be the final size. Then there were fewer than $6\mu^{1/4}$ steps, and both levels next to the
final window have at most $\mu^{3/4}$ vertices outside $P$. We have $k\le\sqrt\mu$, because
$\mu\ge\Theta\ge2^{100}(k+1)^{2}$. By~\eqref{eq:window},
\[
  |\delta X(i,i')|\;\le\;6\mu^{3/4}+3k\bigl(6\mu^{1/4}+3\bigr)\;\le\;33\,\mu^{3/4}
  \;\le\;\frac{\mu}{4\log^{2}\mu}=\frac{\Phi(\mu)}4 .
\]
The last inequality says that $132\log^{2}\mu\le\mu^{1/4}$, and this holds for
$\mu\ge2^{100}$. So the final window contradicts Step~1.
\end{proof}

